\documentclass[11pt]{article}
\usepackage[T1]{fontenc}
\usepackage{lmodern,amsmath,amssymb,amsthm,booktabs}
\usepackage[margin=1in]{geometry}
\usepackage{microtype,xurl}
\usepackage[colorlinks=true,linkcolor=blue,citecolor=blue,urlcolor=blue]{hyperref}
\newtheorem{theorem}{Theorem}
\newtheorem{lemma}[theorem]{Lemma}
\newtheorem{proposition}[theorem]{Proposition}
\newtheorem{corollary}[theorem]{Corollary}
\newcommand{\Z}{\mathbb Z}
\newcommand{\Q}{\mathbb Q}
\newcommand{\F}{\mathbb F}
\newcommand{\OO}{\mathcal O}
\newcommand{\Nm}{\operatorname{N}}
\title{A uniform interval bound for inert-prime mass}
\author{Research note prepared with Codex}
\date{26 October 2026}
\hypersetup{pdftitle={A uniform interval bound for inert-prime mass: quadratic-intervals-v0.1.0},
pdfauthor={Research note prepared with Codex},
pdfsubject={Quadratic characters, divided derivatives, and limits of aggregate prime profiles}}
\begin{document}
\maketitle
\begin{center}\small
Research draft: \texttt{quadratic-intervals-v0.1.0}.\\
Independently cross-checked by another model;
see Section~\ref{sec:checks} for scope and subsequent clarifications.
\end{center}

\begin{abstract}
We give an elementary determinant proof of a uniform bound on the
weighted mass of inert primes in arbitrary intervals. For every
primitive nonprincipal real character $\chi$ of conductor $q$, we prove
\[
 \sum_{\substack{Y<p\le X\\\chi(p)=-1}}\frac{\log p}{p}
 \le \frac12\log\frac XY+\frac14\log q
       +\log(\log q+33)+18\qquad(1\le Y\le X).
\]
For a fixed quadratic field, the allowance above $\tfrac12\log(X/Y)$
depends only on the field, independently of both endpoints. Thus the
same estimate controls every interval $(Y,AY]$ of a fixed ratio $A>1$,
however large $Y$ becomes.

The proof combines divided derivatives with factorial divisibility at
split primes. Their contributions cancel a loss that would otherwise
grow with $\log(X/Y)$, producing an explicit error independent of the
interval endpoints. As a consequence, we bound the logarithmic
derivative of the associated Dirichlet $L$-function:
$L'(\sigma,\chi)/L(\sigma,\chi)
\le\frac12\log q+3\log(\log q+33)+45$ for every real $\sigma\ge1$.
Finally, an explicit continuous model shows that the leading interval
inequalities, together with specified positivity and counting
constraints, are compatible with a residue of order $1/\log q$.
This identifies a limitation of those aggregate constraints alone.
\end{abstract}

\section{Introduction}
Motivated by the broader aims of Tao's challenge \cite{tao}, we turn
from class number one to all quadratic characters.
Let $\Delta$ be a fundamental discriminant, $q=|\Delta|\ge3$, and
$\chi=\chi_\Delta$ its primitive real character. Both signs of
$\Delta$ are allowed. Write $t=\log q$ and
\[
 I_\chi(Y,X)=\sum_{\substack{Y<p\le X\\\chi(p)=-1}}\frac{\log p}{p},
 \qquad C_\chi(X)=\sum_{p\le X}\chi(p)\frac{\log p}{p}.
\]
All logarithms are natural.

\begin{theorem}\label{thm:interval}
For $1\le Y\le X$,
\begin{equation}\label{eq:main}
 I_\chi(Y,X)\le\frac12\log\frac XY+\frac t4+\log(t+33)+18.
\end{equation}
\end{theorem}

Fix a quadratic field and a ratio $A>1$. Theorem~\ref{thm:interval}
gives the same upper bound for the weighted sum of inert primes in
$(Y,AY]$, however large $Y$ becomes. The allowance above
$\tfrac12\log A$ depends only on the field. Thus the estimate places
a uniform restriction on the concentration of inert primes across
all multiplicative scales. For narrow intervals the error may dominate;
the restriction becomes informative for sufficiently wide intervals.

\subsection{Classical uniformity and the explicit refinement}
Endpoint independence for a fixed field already follows from classical
analytic number theory. The prime number theorem in arithmetic
progressions with its classical error term
\cite[Chapter 11]{kedlaya}, followed by partial summation, gives
\[
 I_\chi(1,x)=\tfrac12\log x+c_\chi+o_\chi(1)
 \qquad(x\to\infty).
\]
The remainder is bounded for $x\ge1$. Subtracting the expressions at
$X$ and $Y$ therefore gives, uniformly for $1\le Y\le X$,
\[
 I_\chi(Y,X)=\tfrac12\log(X/Y)+O_\chi(1).
\]
The hidden constant may depend on the field. Theorem~\ref{thm:interval}
supplies the explicit allowance $t/4+\log(t+33)+18$ uniformly as the
field varies, by an elementary determinant argument. Its contribution
is this quantitative refinement within the determinant method;
qualitative endpoint independence is not claimed as new.

The original polynomial argument is due to OpenAI \cite{openai}.
Green's derivative determinant \cite{green} provides the local row
operation. Van Dijk \cite{vandijk} proves nonvanishing by reduction at
an inert prime and gives the cumulative estimate
$I_\chi(1,X)\le\frac12\log X+\frac14\log q+
\log\log(3X)+O(1)$. Kominers \cite{kominers} extends cumulative
prime-mass bounds to every Frobenius class in a Galois extension.
Theorem~\ref{thm:interval} is narrower in its field scope, but controls
an arbitrary interval with an error independent of its location and
length. It is not obtained by subtracting two upper bounds.

Our v7 proof \cite{v7}, Section 4.3, retains factorial divisibility at
split primes to strengthen the inert-prime comparison. Here those
factors have a structural role: they remove a coefficient loss that
would otherwise grow with $\log(X/Y)$. Divided derivatives and their
factorial valuation cost are classical. We use them with the number
of nodes proportional to $R/Y$, where $R$ is the last inert prime
in the interval. The refinement studied here is this combination and
its explicit interval consequence; no novelty is asserted for the
individual ingredients. The source comparison does not establish
priority for the exact numerical bound over all existing literature.

\subsection{Analytic comparisons and consequences}
The associated Dirichlet $L$-function is initially defined by
\[
 L(s,\chi)=\sum_{n=1}^{\infty}\frac{\chi(n)}{n^s},
 \qquad \operatorname{Re}s>1.
\]
Here $L'$ denotes differentiation with respect to $s$. Its logarithmic
derivative $L'/L$ measures the rate of change of $L$ relative to its
value; for real $\sigma>1$ it equals
$\frac{d}{d\sigma}\log L(\sigma,\chi)$.

The quasi-Riemann hypothesis (QRH) supplies stronger asymptotic control.
Taking a zero-free half-plane with fixed boundary $1/2\le\alpha<1$
as input---for example, the
OpenAI result with boundary $\alpha=11/12$ \cite[Theorem 1.1]{qrh}---the
standard explicit formula gives
\[
 \vartheta_\chi(x):=\sum_{p\le x}\chi(p)\log p
 \ll_\alpha x^\alpha\log^2(qx),\qquad x\ge2,
\]
uniformly in $q$ and $\chi$; see \cite[Section 11.1]{kedlaya} for the
explicit formula and the uniform zero count. Here all nontrivial zeros
have real part at most $\alpha<1$; the functional equation bounds
their real parts below by $1-\alpha$, and removing prime powers changes
the estimate by $O(\sqrt x\log^2x)$.
For clarity, the resulting interval comparison is
\[
 I_\chi(Y,X)\le\tfrac12\log(X/Y)
                 +O_\alpha\bigl(\log\log(3q)\bigr)
 \qquad(1\le Y\le X).
\]
Indeed, put $U=(\log(3q))^{3/(1-\alpha)}$. Mertens' estimate bounds
the absolute value of $C_\chi(\min(X,U))$ by $O_\alpha(\log\log(3q))$.
For $X\ge U$, partial summation of the displayed estimate for
$\vartheta_\chi$ gives
\[
 |C_\chi(X)-C_\chi(U)|
 \ll_\alpha U^{\alpha-1}\log^2(qU)=O_\alpha(1).
\]
Thus $C_\chi(X)=O_\alpha(\log\log(3q))$ uniformly in $X$.
Combining the difference $C_\chi(X)-C_\chi(Y)$ with Mertens' estimate,
and discarding the nonnegative ramified-prime term, proves the
comparison. This is a deduction from the cited analytic input.
Its dependence on $q$ is asymptotically stronger than
Theorem~\ref{thm:interval}; we do not compute a numerical constant for
it here. The value of the present argument is its explicit estimate
and elementary mechanism, rather than a stronger asymptotic
prime-distribution theorem.

Cairn \cite{cairn}, Appendix C, already treats all quadratic
characters and transfers the determinant profile to the residue and
regular part of $\zeta(s)L(s,\chi)$. In particular, that paper proves
\[
 \frac{L'(1,\chi)}{L(1,\chi)}<\frac t2+8\log t+45\qquad(t\ge10).
\]
Our interval bound gives the following replacement, without importing
the ideal count or a preliminary lower bound for $L(1,\chi)$.

\begin{corollary}\label{cor:derivative}
Put $B_q=t/2+3\log(t+33)+45$. For $1\le Y\le X$,
\begin{equation}\label{eq:drawdown}
 C_\chi(X)-C_\chi(Y)\ge-B_q.
\end{equation}
For every real $\sigma\ge1$, $L(\sigma,\chi)>0$ and
\begin{equation}\label{eq:derivative}
 \frac{L'(\sigma,\chi)}{L(\sigma,\chi)}\le B_q.
\end{equation}
\end{corollary}

These statements do not assert a new class-number classification,
a fixed zero-free half-plane, or an improvement on the strongest
available bounds obtained using that heavier input. The proof below
uses elementary prime estimates, quadratic arithmetic, and finite
determinants. The $L$-function corollary additionally uses the Euler
product and convergence of the Dirichlet series and its derivative
on the real line at 1. No zero-free theorem is an input.

\section{Elementary prime estimates}
Put $c=2\log2<7/5$, $M(x)=\sum_{p\le x}(\log p)/p$, and
$\vartheta(x)=\sum_{p\le x}\log p$.
We will use, for real $x\ge1$,
\begin{equation}\label{eq:mertens}
 \vartheta(x)\le cx,\qquad \log x-4\le M(x)\le\log x+c.
\end{equation}
Consequently, for $1\le a\le b$ and $a>1$ in the last inequality,
\begin{align}
 \left|M(b)-M(a)-\log(b/a)\right|&<6,\label{eq:interval-mertens}\\
 \sum_{p>a}\frac{\log p}{p^2}&\le\frac{2c}{a},&
 \sum_{p>a}\frac{\log p}{p(p-1)}&\le\frac{2c}{a-1}.
 \label{eq:tails}
\end{align}

For completeness, the usual binomial-coefficient induction gives
$\vartheta(n)\le2n\log2$: for $n=2m$ use $\binom{2m}{m}$ and
the primes in $(m,2m]$; for $n=2m+1$ use
$\binom{2m+1}{m}\le2^{2m}$ and $(m+1,2m+1]$.
The factorial identity gives
\[
 nM(n)-\vartheta(n)\le\log(n!)
 \le n\sum_{p\le n}\frac{\log p}{p-1},\qquad
 \log(n!)\ge n\log n-n.
\]
Partial summation of $\vartheta(x)\le cx$ gives
\eqref{eq:tails}. Retaining the boundary term at 2 also gives
\[
 \sum_p\frac{\log p}{p(p-1)}
 \le\frac{\log2}{2}+\frac32
       \left(c-\frac{\log2}{4}\right)=\frac{25}{8}\log2.
\]
Thus $M(n)\ge\log n-1-(25/8)\log2$. Taking $n=\lfloor x\rfloor$
and using $1+(33/8)\log2<4$ proves the lower bound in
\eqref{eq:mertens}; the upper bound follows from
$\log(n!)\le n\log n$. The range $1\le x<2$ is immediate.
For example, $\log2<7/10$ follows by retaining the first four
terms of the exponential series for $e^{7/10}$.

\section{A determinant adapted to an interval}
Fix for now an integer $K\ge32$. If the interval contains an inert
prime, let $R$ be its largest one. Put
\begin{equation}\label{eq:parameters}
 Z=R/Y>1,\qquad n=\lceil Z\rceil,\qquad
 S=\lfloor n/K\rfloor,\qquad J=\lceil RS/n\rceil.
\end{equation}
We first assume
\begin{equation}\label{eq:nontrivial-range}
 n\ge2K,\qquad n\ge\sqrt q.
\end{equation}
Then $n\le R$, $S\ge2$, $R\ge64$, and $J<R$.

Choose exactly $n$ points $z=u+v\sqrt\Delta$ from the grid
\[
 0\le u<\lceil\sqrt n q^{1/4}\rceil,
 \qquad 0\le v<\lceil\sqrt n q^{-1/4}\rceil,
\]
and put $w=\bar z$. There are at least $n$ grid points, and both
embeddings have modulus less than
\begin{equation}\label{eq:height}
 H=2\sqrt n q^{1/4},\qquad
 \log H=\log2+\tfrac12\log n+\tfrac14t.
\end{equation}
The coordinate ranges have lengths at most $R$, since
$q^{1/4}\sqrt n\le n\le R$. As $R$ is odd and inert,
$\OO_{\Q(\sqrt\Delta)}/R\OO\simeq\F_{R^2}$ and
$1,\sqrt\Delta$ are an $\F_R$-basis. The selected points remain
distinct after reduction.

Use rows $z^rw^s$, $0\le r<R$, $0\le s<S$, and at each point
use the $J$ divided derivatives
\[
 D_z^{[k]}(z^rw^s)=\binom rk z^{r-k}w^s,\qquad 0\le k<J,
\]
with value zero for $k>r$. The entries are algebraic integers.

\begin{lemma}\label{lem:rank}
This matrix has full row rank $N=RS$.
\end{lemma}
\begin{proof}
Modulo $R$, Frobenius gives $w=z^R$. A nonzero row combination
therefore becomes a nonzero polynomial with distinct exponents
$r+Rs$ and degree at most $RS-1$. Its derivatives of orders
$k<J<R$ agree with the prescribed derivatives after substitution;
equivalently, multiply the divided derivatives by the invertible
$k!$. Vanishing of all the columns would give $n$ distinct roots,
each of multiplicity at least $J$, whereas $nJ\ge RS$.
This is van Dijk's argument \cite{vandijk}, with the different
node and derivative counts in \eqref{eq:parameters}.
\end{proof}

Select a nonzero $N$-square minor $d$ and put
$\mathcal D=|\Nm_{\Q(\sqrt\Delta)/\Q}d|\in\Z_{>0}$.
Define the total row degrees and the selected column orders by
\[
 U=\frac{N(R-1)}2,\qquad W=\frac{N(S-1)}2,\qquad
 V=\sum_{\text{selected columns}}k,
\]
and set
\begin{equation}\label{eq:ratios}
 b=\frac WU,\qquad v=\frac VU,\qquad d_0=1-b-v.
\end{equation}
The symbol $d_0$ is a real coefficient, not a discriminant or determinant.
The parameter choices imply
\begin{equation}\label{eq:ratio-bounds}
 0\le b\le\frac1K,\qquad 0\le v<\frac3K,\qquad
 d_0>1-\frac4K\ge\frac78.
\end{equation}
Indeed, $(S-1)/(R-1)\le S/R\le1/K$, and
$J-1<RS/n\le R/K$, so $v<2R/[K(R-1)]<3/K$.

\subsection{Size at both embeddings}
For every nonzero determinant term, the sum of the exponents
$r+s-k$ is $U+W-V$. The bound $\binom rk\le2^r$ therefore gives
\[
 \log\mathcal D\le2\log(N!)+2U\log2+2(U+W-V)\log H.
\]
Since $N\le R^2$ and $R\ge64$,
$\log(N!)/U\le4\log R/(R-1)<1$. Thus
\begin{equation}\label{eq:arch}
 \frac{\log\mathcal D}{2U}
 \le(1+b-v)\log H+\log2+1.
\end{equation}
The absence of a $J\log R$ term is the reason for dividing the
derivative columns by their factorials. The corresponding arithmetic
cost must still be paid, as follows.

\section{Inert and split primes in the same minor}
Put $Q=R/S$. Then $K\le Q<R$ and $Y<Q$.
Write $I=I_\chi(Q,R)$ and let $P$ be the analogous sum over split
primes in $(Q,R]$.

\subsection{Inert primes and the factorial cost}
At an inert prime $p$, Green's integral unitriangular row operation is
\[
 z^rw^s\longmapsto z^{r-pa}(z^p-w)^aw^s,
 \qquad a=\min(\lfloor r/p\rfloor,S-1-s).
\]
Every ordinary derivative of $z^p-w$ is divisible by $p$ after
evaluation, including its zeroth derivative by Frobenius. The
ordinary derivative minor has norm divisible by $p^{2E_p}$, where
\[
 E_p=\sum_{r<R,s<S}\min(\lfloor r/p\rfloor,S-1-s).
\]
Dividing its selected columns by $k!$ gives exactly $d$, hence
\begin{equation}\label{eq:inert-raw}
 \frac{v_p(\mathcal D)}{2U}
 \ge\frac{E_p}{U}-\frac v{p-1}.
\end{equation}
We used $\sum_k v_p(k!)\le V/(p-1)$, and do not claim that
the ordinary-derivative divisibility survives unchanged.

For every integer $a\ge0$,
$\sum_{s<S}\min(a,S-1-s)\ge Sa-a(a+1)/2$.
Using $r/p-1\le\lfloor r/p\rfloor\le r/p$ gives
\begin{equation}\label{eq:E}
 \frac{E_p}{U}\ge
 \frac1p-\frac2{R-1}-\frac1{2Sp}
              -\frac{2R-1}{6Sp^2}.
\end{equation}

\subsection{Split primes}
At a split prime replace $w^s$ by
$(w)_s=w(w-1)\cdots(w-s+1)$. This is another integral
unitriangular change preserving the same selected minor. Each block
of $p$ consecutive factors vanishes in both components of
$\OO/p\OO\simeq\F_p\times\F_p$. Thus its values belong to
$p^{\lfloor s/p\rfloor}\OO$. Divided differentiation in $z$
preserves this factor and integrality. Consequently
\begin{equation}\label{eq:split}
 \frac{v_p(\mathcal D)}{2U}
 \ge\frac{\sum_{r,s}\lfloor s/p\rfloor}{U}
 \ge\frac bp-\frac2{R-1}.
\end{equation}
The inert and split operations are used at disjoint sets of primes,
so their contributions can be added. This is the factorial mechanism
from \cite[Section 4.3]{v7}.

\subsection{A uniform bound for all discarded terms}
Sum \eqref{eq:inert-raw}--\eqref{eq:split} over $(Q,R]$.
Equations \eqref{eq:mertens}--\eqref{eq:tails} show that the total
error in retaining $(1-v)I+bP$ is at most
\begin{equation}\label{eq:error}
 \frac{2cR}{R-1}+\frac{2c}{3}
 +\frac{\log S+6}{2S}+\frac{6c}{K(Q-1)}<6.
\end{equation}
The four terms respectively bound the floor loss across the disjoint
prime sets, the $p^{-2}$ row loss, the $I/(2S)$ term, and the
difference $1/(p-1)-1/p$ in the column cost. For the third term,
$I\le M(R)-M(Q)<\log S+6$.
For $R\ge64$, $S\ge2$, and $Q,K\ge32$, the four bounds are,
respectively, less than $2.845$, $0.934$, $1.675$, and $0.009$.
The sum is less than 6. All inequalities are uniform in $Y$.

Combining this with \eqref{eq:arch} and $1+\log2+6<8$ gives
\begin{equation}\label{eq:coupled}
 (1-v)I+bP\le(1+b-v)\log H+8.
\end{equation}
Let $A$ be the ramified-prime mass in $(Q,R]$. Then
\[
 I+P=M(R)-M(Q)-A,\qquad A\le t/Q.
\]
Substitution in \eqref{eq:coupled}, and
$M(R)-M(Q)\ge\log S-6$, prove
\begin{equation}\label{eq:coupled-two}
 d_0I\le(1+b-v)\log H-b\log S+6b+bt/Q+8.
\end{equation}
Crucially, $1+b-v=d_0+2b$. The $b\log n$ contributed by the
height in \eqref{eq:height} cancels against $b\log S$, up to
$b\log(n/S)$. Discarding the split-prime term would lose this
cancellation and leave an error proportional to $\log(X/Y)/K$.

\section{Completing the constants}
Dividing \eqref{eq:coupled-two} by $d_0$ gives
\begin{align}
 I\le{}&\tfrac12\log n+\tfrac14t+\log2+\frac8{d_0}\notag\\
 &+\frac b{d_0}
     \left(\frac t2+2\log2+\log(n/S)+6+\frac tQ\right).
 \label{eq:before-restore}
\end{align}
The omitted interval has mass at most
\[
 I_\chi(Y,Q)\le\log(Q/Y)+6\le\log(n/S)+6.
\]
Here $R\le nY$, $Y<Q$, and $Q=R/S$.
Also $S\ge n/(2K)$, $n\le2Z$, and
\[
 \frac b{d_0}\le\frac1{K-4},\qquad
 \frac8{d_0}\le\frac{8K}{K-4},\qquad Q\ge K.
\]
It follows that
\begin{equation}\label{eq:K-bound}
 I_\chi(Y,R)\le\frac12\log Z+\frac t4+\log K
 +\frac{K+2}{2K(K-4)}t+C_K,
\end{equation}
where
\begin{equation}\label{eq:CK}
 C_K=\frac52\log2+6+\frac{8K}{K-4}
       +\frac{2\log2+\log(2K)+6}{K-4}<\frac{35}{2}.
\end{equation}
For $K\ge32$ the two variable terms decrease; the derivative of
the last one has numerator
$(K-4)/K-[3\log2+\log K+6]<0$. At 32 the upper bound
$\log2<7/10$ already gives $C_K<17.308<35/2$.

Choose
\[
 K=\lceil t\rceil+32.
\]
Then $K<t+33$ and
\[
 \frac{K+2}{2K(K-4)}t
 \le\frac{(K-32)(K+2)}{2K(K-4)}<\frac12.
\]
Equations \eqref{eq:K-bound}--\eqref{eq:CK} give
\eqref{eq:main} with $R$ in place of $X$ in the nontrivial range.

If \eqref{eq:nontrivial-range} fails, use
$I_\chi(Y,R)\le\log Z+6$. If $n<\sqrt q$, then
$\frac12\log Z\le t/4$. If $n<2K$, then
$\frac12\log Z<\frac12\log(2K)\le\log K$. Either case gives
the same claimed bound, with room to spare. If there is no inert
prime in the interval, the assertion is immediate. Finally,
$I_\chi(Y,X)=I_\chi(Y,R)$ and $R\le X$, so the right side can
only increase when $R$ is replaced by $X$. This proves
Theorem~\ref{thm:interval}.

\section{The signed profile and the logarithmic derivative}
The total ramified mass satisfies
\begin{equation}\label{eq:ramified}
 \sum_{p\mid q}\frac{\log p}{p}
 \le M(t+1)+\frac{t}{t+1}
 <\log(t+33)+\frac{12}{5}.
\end{equation}
Subtract twice \eqref{eq:main} and this ramified contribution from
$M(X)-M(Y)\ge\log(X/Y)-6$. The result is
\[
 C_\chi(X)-C_\chi(Y)
 \ge-\frac t2-3\log(t+33)-\frac{222}{5}>-B_q,
\]
which proves \eqref{eq:drawdown}.

For $\eta>0$, summation by parts gives
\[
 \sum_p\frac{\chi(p)\log p}{p^{1+\eta}}
 =\eta\int_1^\infty C_\chi(x)x^{-1-\eta}\,dx\ge-B_q.
\]
The higher-prime-power terms in the Euler logarithmic derivative are
nonnegative after summing at each prime:
\[
 \sum_{a\ge2}\frac{\chi(p)^a}{p^{a\sigma}}
 =\frac{\chi(p)^2p^{-2\sigma}}{1-\chi(p)p^{-\sigma}}\ge0.
\]
It follows that $L'(\sigma,\chi)/L(\sigma,\chi)\le B_q$ for
$\sigma>1$. Positivity at 1 need not be imported: integrate from
$1+\varepsilon$ to 2 to obtain
$L(1+\varepsilon,\chi)\ge e^{-B_q}L(2,\chi)>0$.
The bounded partial sums of the nonprincipal periodic character
ensure continuity of $L$ and $L'$ on the real line at 1.
Taking the limit gives $L(1,\chi)>0$ and \eqref{eq:derivative}.

The stronger interval statement also gives, for every $Y\ge1$,
\[
 \sum_{p>Y}\frac{\chi(p)\log p}{p^{1+\eta}}
 \ge-B_qY^{-\eta}.
\]
Indeed, integrate $C_\chi(x)-C_\chi(Y)$ from $Y$ to infinity.
This conclusion concerns the signed tail, not its absolute value.

\begin{corollary}\label{cor:split-gap}
For every $Y\ge1$ there is a split prime in
\[
 Y<p\le e^{53}\sqrt q\,(\log q+33)^4Y.
\]
\end{corollary}
\begin{proof}
The split-prime mass in $(Y,X]$ is at least
\[
 \frac12\log(X/Y)-\frac t4-2\log(t+33)-\frac{132}{5},
\]
by subtracting \eqref{eq:main} and \eqref{eq:ramified} from
\eqref{eq:interval-mertens}. At the stated endpoint this is $1/10>0$.
\end{proof}
This illustrates the uniformity in the starting point. It is an
application of this proof, not a claim to improve on the strongest
available analytic prime-distribution theorems.

For comparison with the general residue bound in
\cite[Appendix C.4, (C13)]{cairn}, importing that paper's Fourier
ideal count and regular-part transfer now gives, for $t\ge10$,
\begin{equation}\label{eq:residue}
 L(1,\chi)>
 \frac{\zeta(2)-21/t^3}{t+8\log t+3\log(t+33)+46}.
\end{equation}
This is a corollary relative to those stated inputs, not an additional
self-contained proof of the ideal count. To check the substitution,
take $X_0=\sqrt q\,t^8$ in (C13), use \eqref{eq:derivative},
and use $\gamma-1<0$. The previous denominator was
$t+16\log t+46$. We do not claim a better limiting coefficient
than Cairn's stronger moment and Burgess refinements.

\section{A limit of the aggregate information}
It is tempting to expect that combining enough scale inequalities
will automatically replace $1/\log q$ by $1/\log\log q$ in the
residue bound. The following explicit model shows why a specified
collection of relaxed constraints cannot establish that improvement.
It is not asserted to arise from actual Dirichlet characters.

Use the scaled variable $u=\log x/t$. Let the total prime-log
density be 1, the inert density be
\[
 i(u)=\frac{1+e^{-2u}}2,
\]
and the split density be $(1-e^{-2u})/2$. These densities are
nonnegative and sum to 1. On every interval $0\le a\le b$,
\begin{equation}\label{eq:model-interval}
 \int_a^b i(u)\,du
 =\frac{b-a}{2}+\frac{e^{-2a}-e^{-2b}}4
 \le\frac{b-a}{2}+\frac14.
\end{equation}
Thus the leading part of Theorem~\ref{thm:interval} holds on
\emph{every} interval, not only at a single cutoff.

To retain multiplicative information as well, define the positive
primitive-prime logarithmic measure
\[
 d\nu(u)=\frac{1-e^{-2u}}u\,du\qquad(u>0).
\]
For $z>0$, the elementary identity
\[
 \int_0^\infty\frac{e^{-zu}-e^{-(z+2)u}}u\,du
 =\int_z^{z+2}\frac{dv}{v}=\log\frac{z+2}{z}
\]
shows that its convolution exponential has Laplace transform
\begin{equation}\label{eq:model-euler}
 F(z)=\exp\left(\int_0^\infty e^{-zu}\,d\nu(u)\right)
 =1+\frac2z.
\end{equation}
Equivalently, the resulting primitive-ideal measure is
$\mu=\delta_0+2\,du$. It has the identity ideal, positivity, and an
exactly constant tail density at every positive scale.

We specify the counting constraints retained in this relaxation.
The arithmetic quantity being modeled is the weighted measure
\[
 \mu_\chi=\sum_{n\ge1}\frac{b_\chi(n)}n\,
                \delta_{\log n/t},\qquad
 \frac{\zeta(s)L(s,\chi)}{\zeta(2s)}
       =\sum_{n\ge1}\frac{b_\chi(n)}{n^s},
\]
where $b_\chi(n)\ge0$ counts primitive ideals and $b_\chi(1)=1$.
For a threshold $a\in\{1/2,1/4\}$ and a normalized residue $r>0$
corresponding to $tL(1,\chi)/\zeta(2)$, the retained conditions on
a locally finite measure $\mu$ on $[0,\infty)$ are
\begin{align}
 \mu-\delta_0&\ge0,\label{eq:model-positive}\\
 \mu((u,v])&=r(v-u)\quad(a<u<v),\label{eq:model-count-tail}\\
 1\le\mu([0,a])&\le r(a+\tfrac12).\label{eq:model-count-head}
\end{align}
Here \eqref{eq:model-count-tail} retains the leading constant density
beyond the counting threshold: $a=1/2$ for the finite-Fourier count,
and $a=1/4$ when the Burgess count is imported
\cite[Appendices C.5--C.6, (C15) and (C22)]{cairn}.
The lower bound in \eqref{eq:model-count-head} is the identity ideal.
The upper bound retains the regular-part constraint from
$L'(1,\chi)/L(1,\chi)\le t/2+o(t)$. Indeed, as $z\downarrow0$,
\[
 \int_{[0,\infty)}e^{-zu}\,d\mu(u)
 =\frac rz+\mu([0,a])-ra+O(z),
\]
so the retained bound on the constant term is
$\mu([0,a])-ra\le r/2$.
These are limiting constraints; no finite-$q$ counting error term
or realization by a sequence of characters is imposed.

For example, these conditions imply the first-moment bound
\[
 \int_{[0,V]}u\,d\mu(u)
 \le a\bigl(\mu([0,a])-1\bigr)
       +\frac r2(V^2-a^2),\qquad V\ge a,
\]
by removing the identity atom and using $u\le a$ on the head.
Our model satisfies \eqref{eq:model-positive}--\eqref{eq:model-count-head}
with $r=2$ for both thresholds, since
\[
 \mu([0,a])=1+2a=2(a+\tfrac12),\qquad
 \mu((u,v])=2(v-u)\quad(0\le u<v).
\]
Its first moment is $V^2$, while the displayed upper bound is
$V^2+a^2$. Thus the counting conditions and their first-moment
consequence hold explicitly.

In the usual residue scaling,
$\zeta(1+z/t)\sim t/z$ and $\zeta(2+2z/t)\to\zeta(2)$.
Consequently the model \eqref{eq:model-euler} corresponds to
\[
 \frac{t}{\zeta(2)}L(1+z/t)\longrightarrow z+2,
 \qquad tL(1)\longrightarrow2\zeta(2),
 \qquad \frac1t\frac{L'(1)}{L(1)}\longrightarrow\frac12.
\]
The scaled last expression saturates the leading logarithmic
derivative coefficient. Hence positivity of this convolution
exponential, the constant primitive-ideal tail, and all the leading
interval inequalities are compatible with a residue of order
$1/\log q$.

This is a countermodel only to that limiting relaxation. It omits
the existence of a single periodic character with the prescribed
prime behavior, its functional equation, individual ideal-class
geometry, and other arithmetic correlations. It proves neither
optimality of the determinant method in general nor a limitation
of every possible use of these ideas. It identifies the kind of
additional information a larger improvement would have to retain.

\section{Checks and remaining work}\label{sec:checks}
The accompanying Python program checks exact quadratic-ring minors,
finite-field ranks, inert-prime divisibility with the column-factorial
cost, split-prime row divisibility, the finite row-sum inequality, and
the rational constant bounds. Its small examples deliberately use
more varied derivative-to-node ratios than the final conservative
parameter choice. They check the algebraic mechanisms; they do not
replace the proof over the infinite range. The saved output records
the executed checks and their scope.
There are 25 nonzero-minor checks, 85 inert-prime and 51 split-prime
valuation checks, 55,091 finite row-sum bounds, and 2,286 separately
expanded row-sum identities. In 13 cases, omitting the column-factorial
cost would give a false divisibility claim; the corrected bound passes.

An independent model cross-check, reported on 26 October 2026,
reviewed the ten-page draft and its source and found no substantive
gap in Theorem~\ref{thm:interval}, Corollary~\ref{cor:derivative}, or
Corollary~\ref{cor:split-gap}. The checker reported reproducing all
the supplied checks and verifying another 11 determinant norms through
a separate implementation using ordinary integer block matrices,
covering both real and imaginary fields. It also confirmed the
substitution giving \eqref{eq:residue}, subject to its stated external
ideal-count input. The report is recorded in the accompanying
\texttt{cross\_check.txt}. The two presentation corrections it
requested---the v7 section references and the explicit specification
of the model's counting constraints---are incorporated in this
revision. This is a model cross-check, not formal verification or
human peer review, and it does not establish priority.

The source comparison covers classical analytic estimates, the QRH
comparison, Green, van Dijk, Kominers, the relevant Cairn manuscript,
and the v7 proof; it does not establish priority over all current
manuscripts. The generalization to other Frobenius
classes is not claimed here.

\clearpage
\begingroup\small
\begin{thebibliography}{9}
\bibitem{openai} OpenAI, \emph{Uniform exclusion of Landau--Siegel zeros},
1 October 2026.
\href{https://github.com/openai/math/tree/main/preprints/Uniform-exclusion-of-Landau-Siegel-zeros-October-1-2026}{OpenAI Math Release}.
\bibitem{green} B.~Green, \emph{On Landau--Siegel zeros}, v2, 2026.
\href{https://hexagonmath.org/2610.00311v2}{hexagon:2610.00311v2}.
\bibitem{vandijk} R.~van Dijk, \emph{A Frobenius determinant proof of the
non-existence of Landau--Siegel zeros}, v2, 2026.
\href{https://hexagonmath.org/2610.00409v2}{hexagon:2610.00409v2}.
\bibitem{kominers} S.~D.~Kominers, \emph{Small algebraic integers and
Frobenius prime sums}, v1, 2026, Theorem 1 and Corollary 2.
\href{https://hexagonmath.org/2610.00451v1}{hexagon:2610.00451v1}.
\bibitem{v7} Research note prepared with Codex,
\emph{A unified Frobenius determinant proof for the tenth discriminant},
v0.7.0, 26 October 2026, Sections 3--6.
\href{https://github.com/twinheterodox/tenth-discriminant/releases/tag/v0.7.0}{Release and reproducibility files}.
\bibitem{cairn} M.~Cairn, \emph{Uniform quadratic $L$-value bounds from
Green's determinant and ideal counting}, v3, 2026, Appendix C.
\href{https://rstation.ai/problems/tenth_discriminant/papers/1}{rStation paper 1}.
\bibitem{tao} T.~Tao, \emph{An experimental AI-math challenge: using
tenth discriminant bounds to explore arguments ruling out Siegel zeroes},
20 October 2026.
\href{https://terrytao.wordpress.com/2026/10/20/an-experimental-ai-math-challenge-using-tenth-discriminant-bounds-to-explore-arguments-ruling-out-siegel-zeroes/}{What's New}.
\bibitem{kedlaya} K.~S.~Kedlaya, \emph{Error bounds for primes in
arithmetic progressions}, Chapter 11 of online lecture notes on analytic
number theory, Sections 11.1--11.2.
\href{https://kskedlaya.org/ant/chapter-11.html}{Online lecture notes}.
\bibitem{qrh} OpenAI, \emph{The Quasi-Riemann Hypothesis},
5 October 2026, Theorem 1.1.
\href{https://github.com/openai/math/tree/main/preprints/The-Quasi-Riemann-Hypothesis-October-5-2026}{OpenAI Math Release}.
\end{thebibliography}
\endgroup
\end{document}
